\section{Resumen del Flujo Completo}

La Tabla~\ref{tab:resumen_flujo} condensa el recorrido completo del argumento \texttt{n\_rays\_used:=5} a través de las cuatro capas del sistema.

\begin{table}[H]
  \centering
  \caption{Flujo completo del argumento \texttt{n\_rays\_used:=5} desde la terminal hasta el módulo.}
  \label{tab:resumen_flujo}
  \begin{tabular}{p{1.6cm}p{1.8cm}p{2.2cm}p{1.8cm}}
    \toprule
    Capa & Forma del valor & Mecanismo & Responsable \\
    \midrule
    Terminal &
      string \texttt{'5'} &
      \texttt{sys.argv} / \texttt{split(':=')} &
      OS / ros2 \\[4pt]
    Launch file &
      string \texttt{'5'} $\to$ int \texttt{5} &
      \texttt{.perform(ctx)} + \texttt{int(...)} &
      Usuario \\[4pt]
    Proceso &
      YAML \texttt{n\_rays: 5} &
      \texttt{--ros-args -p} &
      ROS~2 \\[4pt]
    Nodo &
      int \texttt{5} &
      \texttt{get\_parameter .value} &
      rclpy \\[4pt]
    Módulo &
      int \texttt{5} &
      argumento constructor &
      Python \\
    \bottomrule
  \end{tabular}
\end{table}

El diagrama siguiente ilustra las entidades involucradas en cada nivel y los mecanismos de cruce entre ellos:

\begin{figure}[H]
\centering
\begin{lstlisting}[style=bash, caption={Diagrama de flujo Terminal $\to$ Módulo.}, label={lst:diagrama}]
TERMINAL
  "n_rays_used:=5 vmax:=0.8"
      |  sys.argv + split(':=')
      v
LAUNCH FILE  [generate_launch_description]
  DeclareLaunchArgument   <- registra y almacena
  OpaqueFunction          <- logica Python
  .perform(context)       <- resuelve el string
  int('5') = 5            <- conversion: usuario
  Node(parameters={...})  <- descriptor (no proceso)
      |  serializa a YAML, pasa como --ros-args -p
      v
PROCESO  [navigate_individual_pva --ros-args ...]
  rclpy.init()            <- inicia DDS
  Node.__init__()         <- registra en el grafo
  declare_parameter()     <- almacena en param server
  get_parameter().value   <- recupera como int
      |  argumento de constructor Python normal
      v
MODULO  [class PVA]
  self.n_rays_used = 5    <- Python puro
\end{lstlisting}
\end{figure}

\subsection{Observaciones finales}

Tres conclusiones emergen del análisis de este flujo.

\textbf{Todo empieza como string.} La restricción del OS de que \texttt{sys.argv} es siempre una lista de strings se propaga hacia arriba. ROS~2 nunca convierte tipos automáticamente. Esa responsabilidad recae en el usuario dentro del launch file.

\textbf{Cada capa tiene su propio mecanismo de serialización.} El cruce entre capas requiere una conversión de formato: string splitting en la terminal, YAML al pasar al proceso, deserialización en el Parameter Server. Comprender estos mecanismos permite depurar errores de tipo que de otro modo son opacos.

\textbf{El módulo de lógica no sabe que existe ROS~2.} La arquitectura de separación nodo-módulo no es accidental: confina la complejidad del middleware a una capa de adaptación y deja el código de control matemáticamente puro, testeable y portable.
